\sgaStableTarget{sga6.e14.eq.3-1}@@M0@@
where @@M1@@ is a universal polynomial in the @@M2@@, @@M3@@, a polynomial whose coefficients are a priori rational. In fact, one finds [RRR] that these polynomials have \uline{integer} coefficients. This raises the question of the validity of a Riemann--Roch formula ``without denominators''. In fact, such a formula for @@M4@@ cannot in general be sought in @@M5@@ itself, since the homomorphism @@M6@@ is not itself compatible with the @@M7@@-filtrations up to a shift by @@M8@@ (4.6), but is so only up to torsion; thus one cannot define a homomorphism @@M9@@ giving the natural homomorphism @@M10@@ and allowing formula (3.1) to be interpreted as a formula in @@M11@@, and not merely in @@M12@@. But when @@M13@@ is assumed regular, so that @@M14@@ (IV 2.5), one can endow @@M15@@ with the \uline{topological filtration} by the codimension of supports, inherited from the corresponding filtration of @@M16@@. It is plausible that this filtration is compatible with the ring structure, a fact that can be proved at least [RRR] when @@M17@@ is smooth over a field @@M18@@, using the theory of the Chow ring [4]. When this is so, the associated graded
@@M19@@
is a \uline{graded ring}, and, also assuming @@M20@@ regular, we have done what is needed to obtain a natural homomorphism of degree @@M21@@ of graded groups
@@M22@@
Since the @@M23@@-filtration of @@M24@@ is finer than the topological filtration by codimension just considered (X 1.3.3), one obtains a canonical homomorphism of graded groups (and even of graded rings)
@@M25@@
which is moreover an isomorphism mod torsion (Exp. VIII). This allows us, as in [RRR], to consider a theory of Chern classes on @@M26@@ with values in @@M27@@, and likewise on @@M28@@, and thus to interpret formula (3.1) as a formula in @@M29@@. In this form, the formula is proved in [RRR] when @@M30@@ is of \uline{characteristic zero}, at least when @@M31@@ is smooth over a field (a hypothesis serving only to ensure that the topological filtration of @@M32@@ is compatible with the ring structure). The characteristic-zero hypothesis is used essentially
